\documentclass[10pt,a4paper]{amsart}
\usepackage[margin=19mm]{geometry}
\usepackage{amsmath,amssymb,mathtools}
\usepackage[scr=rsfs,cal=euler]{mathalfa}
\usepackage{xfrac}
\usepackage[unicode,hidelinks,bookmarks=false]{hyperref}
\newcommand{\ie}{i.e.\ }
\newcommand{\cf}{cf.\ }
\newcommand{\resp}{respectively\ }
\newcommand{\Subsection}{\subsection}
\providecommand{\nobreakdash}{\nobreak\hbox{-}}
\setlength{\emergencystretch}{2em}
\multlinegap0pt
\usepackage{amssymb}
\usepackage{mathtools}
\usepackage{xcolor}
\newtheorem{theorem}{Theorem}
\newcommand{\mg}{\mathfrak{g}}
\newcommand{\mh}{\mathfrak{h}}
\newcommand{\ms}{\mathfrak{s}}
\newcommand{\weyl}{\mathcal{W}(\ms^\mathbb{C})}
\title{Symplectic Dirac operators on homogeneous spaces}
\author{Dan Ciubotaru, Gang Liu, Salah Mehdi, Nicolas Prudhon}
\date{Source: arXiv:2609.09111v2 (CC BY 4.0)}
\begin{document}
\maketitle

\noindent\emph{Source and licence.} Symplectic Dirac operators on homogeneous spaces. Version arXiv:2609.09111v2 (CC BY 4.0), distributed by arXiv under the Creative Commons Attribution 4.0 International licence, http://creativecommons.org/licenses/by/4.0/, as recorded in the arXiv licence metadata for this exact version and shown on the abstract page https://arxiv.org/abs/2609.09111v2. Full licence notice: \texttt{licenses/arxiv-2609.09111-CC-BY-4.0.txt}.\par\smallskip

\noindent\textit{Editorial scope.} Counted unit: the article's theorem rigidity (source0.tex lines 1062--1070). The article's complete original proof of it is delivered with it (source0.tex lines 1072--1149, one \texttt{proof} environment of 78 lines). No mathematics is written, repaired or re-derived: the statement and the proof are the article's own lines, in the article's own notation, and no retained line is substituted or re-flowed.  Retained uncounted as the source-local context the proof needs: lines 968--983; lines 999--1016.  Nothing was omitted from any retained span, so the package carries no omission record.  No retained line contains a citation, so the wrapper carries no bibliography and no bibliography entry.  No retained line contains a figure environment, an image-inclusion command or drawing code, so no image is delivered and the package declares no asset dependency.  Editorial formatting only, in addition: an AMS \texttt{amsart} preamble that restates the article's own declarations and macros used by the retained lines, namely the environments \texttt{theorem} on separate counters, together with the standard packages those lines need.  The retained environments print in this document as Theorem 1 in order of appearance, whereas the article prints the same items with the numbering of its own full document.  The complete native source of the article as distributed in the arXiv e-print for this exact version is bundled unchanged as \texttt{sources/209779/source0.tex}.
\begin{theorem}
\label{Dirac-Parthasarthy}
The following assertions are equivalent.
\begin{enumerate}
\item The decomposition $\mg= \mh \oplus \ms$ 
is $J$-symmetric.
\item The symplectic Dirac operator $D$ satisfies 
the following formula
which we call the \emph{symplectic Parthasarathy formula} :
\begin{equation}
[D^+,D^-]= \Omega_{\mg}\otimes 1 + \delta(\Omega_{\mh})-1\otimes \tau(\Omega_{\mh})- 2\Omega_\mh \otimes 1 \,.
\label{partha}
\end{equation}
\end{enumerate}

\end{theorem}

\begin{equation} 
\label{partha-sum}
\begin{split}
[D^+,D^-]= &\,\Omega_{\mg}
\otimes 1 - \Omega_\mh \otimes 1 \\ &
-\sum_j w_j \otimes 
\left( \frac12\sum_{kl}B_\mh([Z_k^+,Z_l^-],w_j)
(Z_k^-Z_l^++Z_l^+Z_k^-)\right) \\ &
-\sum_{m}
Z_m^+\otimes \left(\frac{i}2 \sum_{kl} 
B_\ms([Z_k^+,Z_l^-],Z_m^-)
(Z_k^-Z_l^+ + Z_l^+Z_k^-) \right) \\ &
-\sum_l 
Z_m^-\otimes \left( \frac{i}2\sum_{kl}
B_\ms([Z_k^+,Z_l^-],Z_m^+)
(Z_k^-Z_l^+ + Z_l^+Z_k^-) \right)\,.
\end{split}
\end{equation}

\begin{theorem}\label{rigidity}
If the decomposition $\mg= \mh \oplus \ms$ is not $J$-symmetric, then there are no cubic terms $v^\pm \in \weyl$ such that
the operators 
$$\tilde{D}^\pm = D^\pm + v^\pm$$ satisfy the symplectic Kostant-Parthasarathy formula, i.e., for any $v^\pm \in \mathcal{W}(\mathfrak{s}^\mathbb{C})$:
$$
[\tilde{D}^+,\tilde{D}^-] \neq 
\Omega_{\mg}\otimes 1 + \delta(\Omega_{\mh})-1\otimes \tau(\Omega_{\mh})- 2\Omega_\mh \otimes 1+ 1\otimes [v^+,v^-]\,.
$$
\end{theorem}

\begin{proof}

Assume that the decomposition $\mg= \mh \oplus \ms$ is not $J$-symmetric and assume that there exist   $v^\pm \in \weyl$ such that
$$
[\tilde{D}^+,\tilde{D}^-] = 
\Omega_{\mg}\otimes 1 + \delta(\Omega_{\mh})-1\otimes \tau(\Omega_{\mh})- 2\Omega_\mh \otimes 1+ 1\otimes [v^+,v^-]
$$ 
where $\tilde{D}^\pm = D^\pm + v^\pm$.
According to the  equation \eqref{partha-sum} in Theorem \ref{Dirac-Parthasarthy}, we have
\begin{equation*}
[D^\pm,1\otimes v^\mp]=\mp
\frac{i}2\sum_{m}
Z_m^\pm \otimes \left( \sum_{kl} 
B_\ms([Z_k^+,Z_l^-],Z_m^\mp)
(Z_k^-Z_l^+ + Z_l^+Z_k^-) \right)\,.
\end{equation*}
So for each $m$, the cubic terms $v^\pm$  satisfy
the equation:
\begin{equation} 
\label{cubic-equation}  
[Z_m^\pm,v^\pm] =
i\sum_{kl} 
B_\ms([Z_k^+,Z_l^-],Z_m^\pm)
(Z_k^-Z_l^+ + Z_l^+Z_k^-),
\end{equation}
where the commutator on the left hand side stands for the one in $\weyl$.\\ Remember that $\gamma$ defines a linear map $\ms^\mathbb{C} \to \text{End}  (S(\ms^\mathbb{C}))$ that leads to an  
isomorphism 
$$\eta : \weyl \to S(\ms^\mathbb{C})$$ 
by setting $\eta(v)=\gamma(v)(1)$ that we use to identify $\weyl$ with  $S(\ms^\mathbb{C})$.
It should be understood that under this identification, there are two multiplications on $\weyl$.  For $u, v  \in  \weyl $, 
the Weyl multiplication will be denoted by $uv$ as we did previously, and the commutative multiplication of $S(\ms^\mathbb{C})$  will be denoted by $u\cdot v$.

Then we can check directly that 
\begin{equation} 
\label{weyl-sym}  
Z_k^\pm Z_l^\pm=Z_k^\pm \cdot Z_l^\pm, \ \text{and}  \   Z_k^\pm Z_l^\mp=Z_k^\pm \cdot Z_l^\mp \pm \delta_{kl}.
\end{equation} 
Hence for each $m$, the equation  \eqref{cubic-equation} becomes 
\begin{equation} 
\label{cubic-equation-sym}  
[Z_m^\pm,v^\pm] =
i\sum_{kl} 
B_\ms([Z_k^+,Z_l^-],Z_m^\pm)
(Z_k^-\cdot Z_l^+ + Z_l^+\cdot Z_k^-),
\end{equation} 
where the bracket on the left hand side stands for the commutator in $\weyl$.

Since $S(\ms^\mathbb{C})$  is isomorphic to algebra of polynomials and $Z_m^\pm$ is a basis of 
$\ms^\mathbb{C}$,  based on the relation \eqref{weyl-sym}, 
we deduce from the above equation \eqref{cubic-equation-sym} (for all $m$), $v^+$  necessarily is of  the form $v^+=v_1^+ +v_2^+$, 
where  
$$v_1^+ = \sum_{klm,m\leq k}a_{mkl} Z_m^-\cdot Z_k^-\cdot Z_l^+$$  and 
$$v_2^+ = \sum_{klm,m\leq k}b_{mkl} Z_m^+ \cdot Z_k^+\cdot Z_l^+.$$
Moreover, we can check directly that $[Z_m^+,v^+] =[Z_m^+, v_1^+]$. Thus, without loss of generality, we may assume 
$$v^+=v_1^+ = \sum_{klm,m\leq k}a_{mkl} Z_m^-\cdot Z_k^-\cdot Z_l^+.$$
We now have:
\begin{align*}
    [Z_n^+,v^+]
    &= \sum_{m\leq k,l}a_{mkl} [Z_n^+,Z_m^-\cdot Z_k^-\cdot Z_l^+] \\
    &= \sum_{m\leq k,l}a_{mkl} \left(
        [Z_n^+,Z_m^-]\cdot Z_k^-\cdot Z_l^+
        +Z_m^-\cdot [Z_n^+,Z_k^-]\cdot Z_l^+ \right)\\
    &= \sum_{k\geq n,l} 2ia_{nkl}Z_k^-\cdot Z_l^+ 
        + \sum_{m\leq n, l} 2ia_{mnl} Z_m^-\cdot Z_l^+ \\
    &= \sum_{k\geq n,l} 2ia_{nkl}Z_k^-\cdot Z_l^+      +\sum_{k\leq n,l} 2ia_{knl}Z_k^-\cdot Z_l^+\,.
\end{align*}
We observe that each of the $a_{knl}$ terms appears twice: 
once in $[Z_n^+,v^+]$ and another time 
in $[Z_k^+,v^+]$. We then must have
\begin{align*}
     2a_{knl} &= B_\ms([Z_k^+,Z_l^-],Z_n^+)
    =-B_\ms(Z_l^-,[Z_k^+,Z_n^+]) \\
    2a_{knl} &=
    B_\ms([Z_n^+,Z_l^-],Z_k^+)
    =-B_\ms(Z_l^-,[Z_n^+,Z_k^+]),
\end{align*}
which implies that $a_{knl}=-a_{knl}$. So $a_{knl}=0$ for all $a_{knl}$. Hence $v^+=0$. A similar treatment for $v^-$ also shows  that $v^-=0$. This is a contradiction.
\end{proof}

\end{document}
